\section*{Capítulo \ref{ch:b3:lab-techniques-3} --- Técnicas de laboratorio III: la práctica de la investigación}

\begin{solution}{exo:b3:lab-techniques-3:1}
$(0.50/1013)^3 = \num{1.2e-10}$. Las fugas, el gas adsorbido en el vidrio y la desgasificación de la grasa y de los septos
dejan mucho más que eso; un material de vidrio seco y caliente y unas buenas juntas importan más que
ciclos adicionales.
\end{solution}

\begin{solution}{exo:b3:lab-techniques-3:2}
$\qty{20.0}{mmol}/\qty{1.48}{mol/L} = \qty{13.5}{mL}$.
\end{solution}

\begin{solution}{exo:b3:lab-techniques-3:3}
$120 + 10 \times 1.00782503 + 55.93493633 - 0.00054858 = \qty{186.01264}{u}$. Error: $10^6 \times (186.0129 - 186.01264)/186.01264 = +1.4$ ppm,
holgadamente dentro de 5 ppm.
\end{solution}

\begin{solution}{exo:b3:lab-techniques-3:4}
Primaria: el artículo de revista, la patente, la tesis. Secundaria: la revisión, el
manual de constantes. Terciaria: el libro de texto.
\end{solution}

\begin{solution}{exo:b3:lab-techniques-3:5}
$M = \qty{194.0}{g/mol}$: C $96.0/194.0 = 49.48$\,\%, H $10.0/194.0 = 5.15$\,\%, N $56.0/194.0 = 28.87$\,\%. Diferencias de 0.17, 0.07 y 0.16
puntos: dentro de 0.4, el análisis respalda la fórmula.
\end{solution}

\begin{solution}{exo:b3:lab-techniques-3:6}
$\qty{170.0}{mg}/\qty{212.0}{g/mol} = \qty{0.802}{mmol}$; $c = \qty{0.802}{mmol}/\qty{0.54}{mL} = \qty{1.5}{mol/L}$ (dos cifras significativas, fijadas por la
lectura del volumen).
\end{solution}

\begin{solution}{exo:b3:lab-techniques-3:7}
Medias de 73.0 y 77.0\,\%; ambas desviaciones típicas valen $\sqrt{10/4} = 1.58$; $s_p = 1.58$. $t = 4.0/(1.58\sqrt{2/5}) = 4.0$, mayor que
2.31: la mejora es significativa al nivel del 95\,\%.
\end{solution}

\begin{solution}{exo:b3:lab-techniques-3:8}
$F = (0.80/0.30)^2 = 7.1 > 5.05$: el primer método es significativamente menos preciso.
\end{solution}

\begin{solution}{exo:b3:lab-techniques-3:9}
Comprar frascos pequeños con inhibidor, fechar cada uno a su recepción y al abrirlo, guardarlos
cerrados, en la oscuridad y en un lugar fresco; comprobar la presencia de peróxidos antes de cualquier destilación o evaporación
y a intervalos fijos tras la apertura; eliminar los frascos caducados, o con resultado
positivo, por la vía de los residuos peligrosos sin concentrarlos.
\end{solution}

\begin{solution}{exo:b3:lab-techniques-3:10}
Peligros: el hidrógeno liberado (incendio, presión), la reactividad del NaH con el agua
y una descomposición exotérmica conocida de la DMF con el hidruro de sodio que puede descontrolarse
al calentar; la DMF es además tóxica para la reproducción. Medidas: sustituir el disolvente
(THF o 2-metiltetrahidrofurano) o la base; si se mantiene, empezar a pequeña escala, a baja
temperatura, añadiendo el NaH en porciones y siguiendo la temperatura, en una vitrina
tras una pantalla, con salida a través de un borboteador y con un baño de enfriamiento preparado; procedimiento
escrito, una segunda persona presente y, en último lugar, los equipos de protección.
\end{solution}

\begin{solution}{exo:b3:lab-techniques-3:11}
Contribuciones relativas: 0.50, 0.50, $0.010/20.0 = 0.05$, $0.010/15.0 = 0.067$ y 0.10\,\%; combinada:
$\sqrt{0.25 + 0.25 + 0.0025 + 0.0044 + 0.010} = 0.72$\,\%. Dominan las dos integrales: una mejor relación señal/ruido y
una corrección cuidadosa de la línea de base y de la fase importan más que la balanza.
\end{solution}

\begin{solution}{exo:b3:lab-techniques-3:12}
$Q = (89.0 - 82.0)/(89.0 - 80.6) = 0.83 > 0.71$: la prueba lo señala. Antes de rechazarlo, buscar en el cuaderno
una causa (una lectura errónea de la bureta, una burbuja de aire, otro lote de valorante); si se
encuentra, rechazarlo e indicar por qué; si no, conservarlo e informar de la prueba, o repetir la
medida.
\end{solution}

\begin{solution}{pb:b3:lab-techniques-3:1}
\textbf{1.} $\ce{FeCl2 + 2NaC5H5 -> Fe(C5H5)2 + 2NaCl}$.

\textbf{2.} El ciclopentadienuro de sodio se destruye con el oxígeno y el agua, y el cloruro de hierro(II)
se oxida a hierro(III) en el aire húmedo.

\textbf{3.} $\qty{5.00}{g}/\qty{126.8}{g/mol} = \qty{0.0394}{mol}$; $\times \qty{185.8}{g/mol} = \qty{7.33}{g}$; rendimiento: $5.86/7.33 = 80$\,\%.

\textbf{4.} $(0.10/1000)^3 = 10^{-12}$ (en teoría).

\textbf{5.} Calentando el dímero, que experimenta una reacción de retro-Diels--Alder, y
destilando el monómero volátil; a temperatura ambiente, el monómero se vuelve a dimerizar
por una reacción de Diels--Alder, por lo que se mantiene en frío y se usa enseguida.

\textbf{6.} Cada porción libera hidrógeno y calor: añadirlo lentamente mantiene bajo control el flujo de gas
y la temperatura.

\textbf{7.} \qty{186.0126}{u}; error de $+1.4$ ppm.

\textbf{8.} C $120.0/185.8 = 64.58$\,\%, H $10.0/185.8 = 5.38$\,\%; encontrados: 64.41 y 5.47; diferencias de $-0.17$ y $+0.09$, dentro de 0.4.

\textbf{9.} Los diez hidrógenos son equivalentes por simetría, y los anillos giran deprisa.

\textbf{10.} Una sola señal, ya que los diez carbonos son equivalentes.

\textbf{11.} Una estructura de rayos X de monocristal.

\textbf{12.} No: como complejo de 18 electrones (\cref{ch:b3:organometallic-bonding}) es estable al aire, a diferencia de sus
precursores.

\textbf{13.} Sus señales no se solapan con las del ferroceno, es un sólido puro, no volátil y
no higroscópico que puede pesarse con exactitud, y es inerte frente a la
muestra.

\textbf{14.} La relajación completa de cada señal entre barridos (una espera de varios $T_1$),
una excitación uniforme, una relación señal/ruido suficiente y una corrección cuidadosa de la fase y de la línea
de base.

\textbf{15.} Ferroceno, $\ce{C10H10Fe}$, \qty{185.8}{g/mol}; 1,3,5-trimetoxibenceno, $\ce{C9H12O3}$, \qty{168.0}{g/mol}.

\textbf{16.} $P_x = 3.942 \times (3/10) \times (185.8/168.0) \times (15.0/20.0) \times 0.999 = 0.980$: 98.0\,\%.

\textbf{17.} Relativa: $\sqrt{0.50^2 + 0.50^2 + 0.05^2 + 0.067^2 + 0.10^2} = 0.72$\,\%; absoluta: $0.72\,\% \times 98.0 = 0.7$ puntos porcentuales.

\textbf{18.} Sí, pero débilmente: una impureza del 2\,\% de composición parecida cambia C y H en menos
de lo que el análisis puede detectar; la qNMR es la medida de pureza más informativa.

\textbf{19.} El hidruro de sodio (hidrógeno con el agua, incendio), el hidrógeno desprendido, el THF y el
ciclopentadieno (muy inflamables; el THF, \omterm{def:b3:lab-techniques-3:peroxide}{formador de peróxidos}), el craqueo en caliente del
dímero, el cloruro de hierro(II) (irritante) y el ferroceno (sólido inflamable, nocivo).

\textbf{20.} Hidruro de sodio: pesado como dispersión en aceite bajo gas inerte, añadido en
porciones a una disolución agitada y enfriada, con salida del gas a través de un borboteador y sin agua
cerca. THF: seco y sin peróxidos, procedente de una columna, usado en la vitrina bajo nitrógeno.

\textbf{21.} Bajo gas inerte y con enfriamiento, por adición lenta de isopropanol, después de
etanol y después de agua, esperando cada vez a que cese el desprendimiento de gas.

\textbf{22.} Una onda reversible de un electrón, $\ce{Fe(C5H5)2+}/\ce{Fe(C5H5)2}$, con picos separados unos 59 mV a \qty{25}{\celsius}
(\cref{ch:b3:electrode-kinetics}); el par es rápido, estable y casi insensible al
disolvente, lo que lo convierte en una referencia interna cómoda.

\textbf{23.} El cuaderno: fecha, esquema, cantidades y procedencias, la \omterm{def:b3:lab-techniques-3:assessment}{evaluación de riesgos}, lo que
se hizo y se observó, los rendimientos y los nombres de los ficheros de \omterm{def:b3:lab-techniques-3:notebook}{datos brutos}. La \omterm{def:b3:lab-techniques-3:literature}{información
complementaria}: el procedimiento completo y todos los datos de caracterización, con copias de los
espectros.

\textbf{24.} Para que otros, y los propios autores años después, puedan volver a abrirlos y procesarlos
sea cual sea el programa del que dispongan: localizables, accesibles, interoperables y reutilizables.

\textbf{25.} La pureza por qNMR de la muestra de ferroceno es de $98.0 \pm 0.7$\,\% (incertidumbre típica).
\end{solution}
